\begin{lemma}[Extension base: technical assertion (iv)]
Let \(T\) be a positive-rung extended Tashkinov tree generated by

\[
 \mathcal T=\{(T_i,\varphi_{i-1},S_{i-1},F_{i-1},\Theta_{i-1}):
                   1\leq i\leq n+1\},
\]

and suppose that \(T\) has the maximum property.  Keep the graph \(G\),
palette size \(k\), and distinguished uncolored edge \(e_0\) fixed, but
allow the certified input coloring to vary arbitrarily.  Assume explicitly
the following six assertions for every positive-rung ETT \(T'\) on these
fixed data, with arbitrary certified input coloring \(\varphi'_0\), maximum
property, and rung number \(s<n\), writing \(\varphi'_s\) for its generating
coloring and \(\mathcal T'\) for its certified series:
\begin{enumerate}
\item \(V(T')\) is elementary with respect to \(\varphi'_s\).
\item Under every proper \((T'_{s+1},D'_s,\varphi'_s)\)-stable coloring
leaving precisely \(e_0\) uncolored, \(T'_{s+1}\) is interchangeable for
any pair of colors of which at least one is missing on \(T'_{s+1}\).
\item For each PE slot \(a\leq s\) with supporting vertex \(u_a\) and
\(m(u_a)=b\), every proper \((T'_a,D'_a,\varphi'_a)\)-stable coloring
leaving precisely \(e_0\) uncolored is
\((T'(u_a)-u_a,D'_{b-1},\varphi'_{b-1})\)-stable.  When \(b>1\), it is
also \((T'_{b-1},D'_{b-1},\varphi'_{b-1})\)-stable.  Moreover, the stored
color pairs at two distinct PE slots with distinct supporting vertices
and the same first-appearance index are disjoint.
\item If the terminal slot is PE with supporting vertex \(u_s\) and
stored colors \(\gamma'_s,\delta'_s\), then under every proper
\((T'_s,D'_s,\varphi'_s)\)-stable coloring leaving precisely \(e_0\)
uncolored, the \((\gamma'_s,\delta'_s)\)-component through \(u_s\)
contains exactly the one vertex \(u_s\) of \(T'_s\).
\item If \(\tau\) is any proper \((T'_s,D'_s,\varphi'_s)\)-stable
coloring leaving precisely \(e_0\) uncolored, \(d\) is defective on
\(T'_s\) under \(\tau\), and \(x\) is a nonminimum vertex of
\(I[\partial_{\tau,d}(T'_s)]\), then
\(x\preceq_{T'}u_a\) for every PE supporting vertex or SE extension
vertex \(u_a\) with \(m(x)\leq a\leq s\).
\item Every proper \((T'_s,D'_s,\varphi'_s)\)-stable coloring leaving
precisely \(e_0\) uncolored is a \(\varphi'_s\bmod T'_s\) coloring.
\end{enumerate}

Let the terminal certified slot of \(\mathcal T\) have kind PE, exact
supporting vertex \(v_n\), and exact gamma and delta fields
\(\gamma_n,\delta_n\).  Then for every proper
\((T_n,D_n,\varphi_n)\)-stable coloring \(\sigma_n\) leaving precisely
\(e_0\) uncolored and every \(w\in V(T_n)\),
\[
 w\in P_{v_n}(\gamma_n,\delta_n,\sigma_n)
       \quad\Longleftrightarrow\quad w=v_n.
\]
Thus the terminal two-color component meets \(T_n\) exactly at its
supporting vertex.
\end{lemma}
\begin{proof}
Fix the terminal PE data and \(\sigma_n\).  The component in
\noderef{TwoColorComponent} contains its root \(v_n\), so the reverse
implication in the displayed equivalence is immediate.  Suppose, for a
contradiction, that it also contains a vertex
\(w\in V(T_n)\setminus\{v_n\}\).  We construct all data needed for the
two cases below from the hypotheses of this lemma.

Apply \noderef{ExtensionRules} to \(T\) and its given maximum property.
That node asks only for lower ETTs whose input coloring is the original
input coloring of \(T\).  Specialize the broader package in this lemma to
\(\varphi'_0=\varphi_0\) and to each such lower ETT; all six conclusions
are unchanged by this specialization.  The conclusion of
\noderef{ExtensionRules} therefore gives current assertion (iii), including
unconditional stability before every current PE support, predecessor-stage
stability when its first-appearance index is greater than one, and the
unconditional disjointness conclusion for equal first-appearance indices.  In particular, if
\[
                         j=m(v_n),
\]
then \(\sigma_n\) is
\((T_j(v_n)-v_n,D_{j-1},\varphi_{j-1})\)-stable.  This invokes no
unstated current-rung hypothesis.

We record the remaining common preparation.  Let \(q\) be the exact
predecessor slot representing \(n\), choose the exact predecessor slot
\(j_{\rm slot}\) representing \(j\), and let
\(T_j(v_n)\) be the \noderef{TreeSequencePrefix} ending at the displayed
occurrence of \(v_n\) in \(T_j\).  The ordered-stage clauses of
\noderef{TashkinovSeriesStep} and the definition of
\noderef{ETTFirstAppearanceIndex} give
\[
 j_{\rm slot}.\mathrm{val}+1=j\leq n=q.\mathrm{val}+1,
 \qquad T_{j-1}\prec T_j(v_n)\preccurlyeq T_j\preccurlyeq T_n,
\]
with all edge identities retained.  For every earlier positive target
rung \(s<n\), \noderef{CertifiedETTRestriction} gives an actual ETT on
the original input \(\varphi_0\), with rung \(s\), displayed tree the
original slot \(T_{s+1}\), generating coloring \(\varphi_s\), and
inherited maximum property.  Its returned pointwise equalities preserve
all retained raw slots and exact connecting-edge identities, its
computed-set equalities preserve every retained \(D_a\), and its
heterogeneous equalities transport all retained step certificates.  Thus
every use of the six displayed hypotheses on such a stage instantiates
\(\varphi'_0=\varphi_0\) explicitly.
At the initial rung, the two initial conclusions used in this preparation
are supplied by \noderef{TechnicalTheoremBase}.

Put \(L=\noderef{ExtensionBaseActiveIndices}(v_n)\) and partition it into
the nonempty finite support classes
\(L(x)=\noderef{ExtensionBaseSupportClass}(v_n,x)\).  Applying lower-rung
assertion (v) to the exact restricted stages verifies the explicit order
premise of \noderef{ExtensionBaseSupportOrder}.  We obtain
\[
 s<t\text{ in }L\Longrightarrow v_t\preceq v_s,
 \qquad m(v_s)=j,
 \qquad v_n\preceq v_s.                                    \tag{1}
\]
The same restricted instances of lower assertions (i) and (v) meet the
two premises of \noderef{ExtensionBaseDeltaExclusion}; hence
\[
 \delta_t\notin\overline\varphi_{s-1}(T_s),\quad
 \delta_t\ne\gamma_s\quad(s\leq t\text{ in }L),             \tag{2}
\]
and every such \(\delta_t\) is absent from
\(\overline\varphi_{j-1}(T_j)\).

Enumerate each class increasingly.  The exact PE fields in
\noderef{TashkinovSeriesStep}, together with (2), satisfy
\noderef{ExtensionBaseSupportClassColors}; consequently adjacent slots
in a class obey \(\gamma_{a+1}=\delta_a\), and
\(\gamma_0,\delta_0,\ldots,\delta_p\) are injectively indexed.  The
full lower package, (1), and (2) satisfy
\noderef{ExtensionBaseStoredPairOverlap}.  Thus two active stored pairs
overlap exactly when their slots are consecutive in the same support
class, and their intersection is then both
\(\{\gamma_t\}\) and \(\{\delta_s\}\).                     \tag{3}

The unsupported-vertex and full-class conclusions of
\noderef{ExtensionBaseClaimMissingColors} give Claim 4.1(a3)--(a4),
including for every class \(i_0<\cdots<i_p\)
\[
 \overline\sigma_n(x)\cap\bigcup_{a=0}^p S_{i_a}
      =\{\delta_{i_p}\},
 \qquad
 \overline\varphi_{j-1}(x)
      =(\overline\sigma_n(x)\setminus\{\delta_{i_p}\})
          \cup\{\gamma_{i_0}\}.                            \tag{4}
\]
Lower assertion (ii), applied to each earlier active PE restriction,
meets the sole extra premise of
\noderef{ExtensionBaseClaimStrongInterchangeable}; every active
nonterminal stored pair is therefore strongly interchangeable under
\(\sigma_n\).  Finally, \noderef{ExtensionBaseClaimTreeEdgeColors}
fixes every exact displayed edge of \(T_j\) under \(\sigma_n\).

Equations (1)--(4) are precisely the four preparation groups in
\noderef{ExtensionBaseAlgorithmA}.  Apply that node to obtain its finite
decreasing restoration trace, and apply \noderef{ExtensionBaseClaim43}
to its invariant and the stability furnished by current assertion (iii).
Writing \(\varrho_1\) for the output, we have the exact conclusions
\begin{enumerate}
\item[(b1)] \(\varrho_1\) is
  \((T_j(v_n)-v_n,D_{j-1},\varphi_{j-1})\)-stable;
\item[(b2)] its missing sets agree with \(\varphi_{j-1}\) on
  \(V(T_j)\setminus\{v_n\}\), agree with those of \(\sigma_n\) at
  \(v_n\), and every displayed edge of \(T_j\) has its
  \(\varphi_{j-1}\)-color;
\item[(b3)] every changed exact edge is outside \(G[V(T_j)]\), and its
  two colors lie in the union of stored pairs belonging to restoration
  classes other than the class of \(v_n\);
\item[(b4)] all preterminal pairs in the class of \(v_n\) remain
  strongly interchangeable under \(\varrho_1\).
\end{enumerate}

We next construct the closure used in both cases.  If \(j>1\), apply
\noderef{CertifiedETTRestriction} with target rung \(r=j-1\); the
inequalities \(1\leq j-1<n\) follow from \(j>1\) and \(j\leq n\).
Call the returned ETT \(B\).  Its terminal ladder stage used by
\noderef{ModColoring} is exactly the old \(T_{j-1}\), its displayed next
tree is exactly \(T_j\), its generating coloring is
\(\varphi_{j-1}\), its input is \(\varphi'_0=\varphi_0\), and it has
maximum property.  The returned data and computed-set equalities identify
all of these expressions with the original series slots.
Restrict (b1) to \(T_{j-1}\); its vertex and controlled-edge clauses are
exactly \noderef{StableColoring} there.  Assertion (vi) from precisely this
original-input instance makes \(\varrho_1\) a \noderef{ModColoring} of
\(B\).  The PE versus SE/RE
constructor of \noderef{TashkinovSeriesStep} supplies the exact eligible
route from \(T_{j-1}\) to \(T_j(v_n)\), including the named connecting
edge in the SE/RE alternative.  To supply the elementarity premise of
\noderef{ExtensionBaseClosure}, take any corresponding certificate \(C\).
Its lower ETT is \(C.\mathrm{replacementETT}\), with certified input
\(C.\mathrm{replacementInput}\), rather than \(\varphi_0\).  The equality
\(C.\mathrm{rungAgreement}\), together with the strict stage bound for
\(B\), puts this replacement ETT strictly below \(n\), and its positive
rung condition is part of its ETT certificate.  First apply
\noderef{MaximumPropertyTransport} to obtain maximum property for
\(C.\mathrm{replacementETT}\); then instantiate the broadened package with
input argument \(C.\mathrm{replacementInput}\) and lower-ETT argument
\(C.\mathrm{replacementETT}\), and extract assertion (i).  Consequently
\noderef{ExtensionBaseClosure} gives the required \(\varrho_1\)-closure
\(T'_j\).

If \(j=1\), use instead the initial certificate in
\noderef{TashkinovSeriesStep}.  It is a \noderef{KTriple}, and its
initial closure displays \(T_1\) as a Tashkinov augmentation sequence
from the distinguished edge.  The b1 missing-set equality and b2
exact-edge equality copy every augmentation through \(T_1(v_n)\) under
\(\varrho_1\).  Properness and unique uncoloredness of the restoration
output, together with the graph and numerical clauses of the initial
\(k\)-triple, make \((G,e_0,\varrho_1)\) another \(k\)-triple.  The
exact prefix through \(v_n\), followed by any chosen closure under
\(\varrho_1\), therefore remains a Tashkinov tree by
\noderef{TashkinovAugmentationSequencePreservesTashkinovTree}.  Choose
such a closure using \noderef{TashkinovClosureExists}; then
\noderef{TashTree} supplies its elementariness.  In this initial case,
the first clause of \noderef{MaximumProperty} identifies
\(|V(T_1)|\) with \noderef{TashkinovOrder}, so any later construction
containing \(T_1\) and one additional vertex is already contradictory.

Thus in either case we have a \(\varrho_1\)-closure \(T'_j\) of
\(T_j(v_n)\), with the relevant maximum certificate when \(j>1\), and
\[
              V(T'_j)\text{ elementary under }\varrho_1.    \tag{5}
\]
The finite missing-color intersection gives the exhaustive split
\[
 \overline\varrho_1(T'_j-v_n)\cap S(L_1)\ne\varnothing
 \quad\text{or}\quad
 \overline\varrho_1(T'_j-v_n)\cap S(L_1)=\varnothing,       \tag{6}
\]
where \(L_1=L(v_n)\).  We treat the two alternatives separately.

In the first alternative, apply
\noderef{ExtensionBaseCaseOneSelection}.  Its exact first-position
choice yields \(u\in V(T'_j)\setminus V(T_j)\), the least missing gamma
at \(u\), and equations (14)--(15) for the processed prefix of \(L_1\).
These conclusions, (3), and (b4) meet the hypotheses of
\noderef{ExtensionBaseAlgorithmB}; let \(\varrho_2\) be the resulting
trace output.  The exact conclusion of \noderef{ExtensionBaseClaim44}
is
\begin{align*}
 &\varrho_2\text{ is }(T_j(v_n)-v_n,D_{j-1},\varphi_{j-1})
       \text{-stable},                                      &&\tag{c1}\\
 &\overline\varrho_2=\overline\varrho_1
       \text{ on }V(T_j)\cup V(T'_j(u)-u),
   \quad \varrho_2=\varrho_1
       \text{ on every displayed edge of }T_j\cup T'_j(u), &&\tag{c2}\\
 &\gamma_{1_1}\in\overline\varrho_2(u).                    &&\tag{c3}
\end{align*}

We now derive the Case 1 contradiction, retaining the exact
augmentations.  When \(j>1\), restrict (c1) to \(T_{j-1}\); assertion
(vi) from the original-input instance \(\varphi'_0=\varphi_0\) makes
\(\varrho_2\) a mod-coloring of \(B\).  Each
augmentation from \(T_{j-1}\) to \(T_j(v_n)\) keeps its old
missing-color witness by (c1)
and keeps its exact new edge color by (b2),(c2).  The same verification
at every retained closure step before \(u\) uses (c2) on the old witness
and on the exact appended edge.  Thus, when \(j>1\), \(T'_j(u)\) carries
a certified corresponding certificate \(C_2\) under \(\varrho_2\).
Here the relevant lower ETT is \(C_2.\mathrm{replacementETT}\) with input
\(C_2.\mathrm{replacementInput}\): rung agreement and the stage bound for
\(B\) give its strict lower-rung inequality, and
\noderef{MaximumPropertyTransport} gives its maximum property.  Instantiating
the broadened package with that replacement input and replacement ETT then
supplies assertion (i) for this exact certificate; separately, the
transported maximum property is the certificate used in the scan below.  When
\(j=1\), the identical edge-by-edge verification begins at the initial
edge and \noderef{TashkinovAugmentationSequencePreservesTashkinovTree}
says that the resulting prefix is a Tashkinov tree under \(\varrho_2\).

The supporting-PE transport in
\noderef{SupportingPEMissingTransport}, (b2), and (c2) give
\[
 \overline\varphi_{j-1}(v_n)
   \subseteq\overline\varrho_2(v_n)\cup\{\gamma_{1_1}\}
   \subseteq\overline\varrho_2(v_n)\cup\overline\varrho_2(u). \tag{7}
\]
Starting with \(T'_j(u)\), scan the displayed positions of \(T_j\) in
order.  Skip a position whose new vertex is already present.  At every
other position, the exact \noderef{TreeSequence} attachment edge joins
the new vertex to an old vertex.  Its \(\varphi_{j-1}\)-color is missing
at an old witness in the original
\noderef{TashkinovAugmentationSequence}.  If that witness is not \(v_n\),
(b2),(c2) preserve its complete missing set; if it is \(v_n\), (7)
supplies the same color at \(v_n\) or at the already present vertex
\(u\).  The exact edge color is preserved by (b2),(c2).  Hence every
nonskipped position has an actual \noderef{TashkinovAugmentation}
certificate.  Induction over this finite scan produces the required
eligible route when \(j>1\), and a Tashkinov tree when \(j=1\), containing
\(V(T_j)\cup V(T'_j(u))\).  Taking its closure with
\noderef{TashkinovClosureExists} preserves this inclusion.  In the first
case the mod-coloring clause supplies a corresponding certificate
\(C_{\rm scan}\) for this closure.  Its certified ETT has input
\(C_{\rm scan}.\mathrm{replacementInput}\), and rung agreement preserves
the strict bound inherited from \(B\).  Apply
\noderef{MaximumPropertyTransport} first, and then instantiate the broadened
package with this replacement input and replacement ETT to extract assertion
(i) for \(C_{\rm scan}\).  The transported maximum property also supplies the
stage bound used next.  Since
\(u\notin V(T_j)\), the closure has strictly more vertices than \(T_j\),
whereas the stage-\(j\) clause of the original
\noderef{MaximumProperty} bounds every such closure by \(T_j\) when
\(j>1\).  When \(j=1\), it is a Tashkinov tree strictly larger than
\(T_1\), contradicting the initial maximum equality with
\noderef{TashkinovOrder}.  This is a contradiction in either case.

It remains to exclude the second alternative in (6).  Equation (4) for
the terminal class and (b2) give
\[
             \overline\varrho_1(v_n)\cap S(L_1)=\{\delta_n\}.
\]
Together with the second alternative and the closure certificate, this
is exactly the premise of \noderef{ExtensionBaseCaseTwoObservation19};
we obtain equation (19), including its exact appended-edge exclusion.
The terminal PE constructor and
\noderef{HkuParallelBoundaryStructure} give the unique
\(\gamma_n\)-boundary edge \(f_n\) and locate every
\(\delta_n\)-boundary edge at the seed.  By (3), the restoration colors
avoid \(\gamma_n,\delta_n\); hence (b3) and terminal stability meet all
hypotheses of \noderef{ExtensionBaseCaseTwoObservation20}.  We obtain
the two exact boundary equalities (20).

We first root the two-color path rather than choosing an unspecified
orientation.  Equation (19) says
\(\delta_n\in\overline\varrho_1(v_n)\).  By
\noderef{MissingColors} and \noderef{PresentColors}, no
\(\delta_n\)-colored identity is incident with \(v_n\).  On the other
hand, (20) says that the exact identity \(f_n\) has
\(\varrho_1\)-color \(\gamma_n\), and the terminal PE data say that it is
incident with \(v_n\).  Any second incident \(\gamma_n\)-identity would
be adjacent to \(f_n\) by \noderef{AdjacentEdges} and would violate
properness in \noderef{IsProperEdgeColoring}.  Thus the vertex-node
\(v_n\) has degree exactly one in
\noderef{TwoColorIncidenceGraph} for
\((\varrho_1;\gamma_n,\delta_n)\).  In particular
\(\gamma_n\ne\delta_n\), since \(f_n\) makes \(\gamma_n\) present at
\(v_n\) while equation (19) makes \(\delta_n\) missing there.

Apply \noderef{TwoColorComponentSpanningPath} to this rooted component.
The resulting full-support simple path contains \(v_n\).  It cannot have
\(v_n\) as an internal node, since that would give \(v_n\) two path
neighbors, contradicting the degree-one calculation.  Reorient the path
so that its first endpoint is \(v_n\), and call its other endpoint \(x\).
The endpoint \(x\) is a vertex-node.  Indeed, every colored edge-node has
the two distinct endpoint vertex-nodes as neighbors by
\noderef{TwoColorIncidenceGraph} and
\noderef{LooplessMultigraph}.  If
such an edge-node were the other endpoint of the spanning path, its unused
second neighbor already lies on the full-support path.  That neighbor
would then acquire a third incidence when internal, or would be \(v_n\)
and give \(v_n\) degree two.  Both alternatives contradict
\noderef{TwoColorIncidenceDegreeAtMostTwo} and the degree-one calculation
at \(v_n\).  Hence \(x\) is a well-defined graph vertex distinct from
\(v_n\).  The b3 edge-control clause, together with the exclusion of
\(\gamma_n,\delta_n\) from the restoration colors, says edge by edge that
the \((\gamma_n,\delta_n)\)-incidence graphs under \(\sigma_n\) and
\(\varrho_1\) coincide.  Thus this is also the full contrary component
specified under \(\sigma_n\).

We now exclude \(x\) from the two tree sets vertex by vertex.  At the
pre-terminal PE coloring, \noderef{ParallelExtensionColor} makes
\(\gamma_n\) missing at \(v_n\); the selected stability transfers this
to \(\varphi_{n-1}\).  If the preceding restriction has positive rung,
lower assertion (i), applied with \(\varphi'_0=\varphi_0\), makes
\(V(T_n)\) elementary under \(\varphi_{n-1}\); in the initial-rung case
the same conclusion is supplied by \noderef{TechnicalTheoremBase}.
Thus \(\gamma_n\) is
present at every \(y\in V(T_n)\setminus\{v_n\}\), by
\noderef{MissingColors} and \noderef{PresentColors}.  Equation (2) with
\(s=t=n\) says that \(\delta_n\) is present at every vertex of \(T_n\)
under that same coloring, using the union in
\noderef{TreeMissingColors} and the same complement definitions.  The PE
selected-stability clauses give the same
two present-color facts at every such vertex under the coloring on which
the terminal component switch is performed.  In that switch, a vertex
incident with both colors has both incident edge-nodes in the switched
component or neither, by \noderef{TwoColorIncidenceGraph} and
\noderef{TwoColorComponent}; therefore both colors remain present there.  The
vertex clause of terminal \noderef{StableColoring} preserves both at
every such \(y\) under \(\sigma_n\), and b3 preserves them under
\(\varrho_1\) because neither terminal color is a restoration color.
Consequently every vertex-node of \(T_n-v_n\) has two neighbors in the
two-color incidence graph.  Since the endpoint \(x\) has degree one, it
cannot be any such vertex; since \(x\ne v_n\), we have
\(x\notin V(T_n)\).

Likewise, if \(x\in V(T'_j)\setminus\{v_n\}\), the second alternative in
(6) says that neither \(\gamma_n\) nor \(\delta_n\), both members of
\(S(L_1)\), is missing at \(x\).  The complement relation in
\noderef{MissingColors} and the witness definition
\noderef{PresentColors} would again give two distinct incidence neighbors,
contrary to the degree-one endpoint.  The remaining possibility
\(x=v_n\) was already excluded, so \(x\notin V(T'_j)\).  Since our
contradiction assumption puts a second vertex of \(T_n\) on this exact
rooted path,
\noderef{ExtensionBaseCaseTwoObservation21} yields an exact exit suffix
ending at
\[
                         u\in V(T'_j)\setminus V(T_n).       \tag{8}
\]

The initial \noderef{KTriple} has \(k\geq\Delta(G)+1\).  By
\noderef{Degree} and \noderef{MaximumDegree}, at most \(\Delta(G)\)
exact edge identities are incident with \(u\).  Distinct colored
identities there are adjacent by \noderef{AdjacentEdges}, so properness
in \noderef{IsProperEdgeColoring} assigns them distinct colors.
Consequently \noderef{PresentColors} has cardinality at most
\(\Delta(G)<k\), and its complement in \([k]\), as defined by
\noderef{MissingColors}, contains some color
\(\beta\in\overline\varrho_1(u)\).  The Case 2 equality in (6) and
\(u\ne v_n\) give \(\beta\notin S(L_1)\), in particular
\(\beta\ne\gamma_{1_1}\).  Elementariness (5), (b2), the supporting-PE
missing transport, and the exact augmentation route now meet the
hypotheses of \noderef{ExtensionBaseCaseTwoObservation23}.  It chooses
the exact \(z\)-position and supplies all vertex and edge exclusions of
equation (23).

Write the exact terminal support class as
\[
                 L_1=\{i_0<i_1<\cdots<i_p=n\}.
\]
When \(p>0\), put \(n^-=i_{p-1}\).  The adjacent-chain conclusion of
\noderef{ExtensionBaseSupportClassColors} gives the exact identity
\(\delta_{n^-}=\gamma_n\) required by
\noderef{ExtensionBaseClaim45}.  We now verify its only additional
ambient coloring premise.  Define
\[
 \varrho_0=\noderef{OutsideTreeColorSwap}
      (T'_j,\varrho_1,\beta,\delta_n),\qquad
 \mu_0=\noderef{TwoColorComponentSwap}
      (G,\varrho_0,\gamma_n,\beta,u).
\]
Every support color is present at \(u\) under \(\varrho_1\), because
\(u\ne v_n\), the Case 2 intersection there is empty, and
\noderef{MissingColors} is the complement of the incidence-witness set
\noderef{PresentColors}.  The outside
swap fixes every exact edge incident with the in-tree vertex \(u\), so
the same colors are present there under \(\varrho_0\), while \(\beta\)
is missing.  Properness, \noderef{AdjacentEdges}, and
\noderef{IsProperEdgeColoring} make the incident \(\gamma_n\)-identity
unique, and \noderef{TwoColorIncidenceGraph} puts its edge-node in the
component rooted at \(u\).  The definition
\noderef{TwoColorComponentSwap} changes that identity to \(\beta\) and
changes only the two displayed colors.  Thus it makes \(\gamma_n\)
missing but cannot remove an incident edge of any preterminal gamma.  Indeed,
\(\beta\notin S(L_1)\), and the injective color chain from
\noderef{ExtensionBaseSupportClassColors} makes every preterminal gamma
distinct from \(\gamma_n\).  Hence every preterminal gamma is present at
\(u\) under \(\mu_0\), exactly as required by the final explicit premise
of \noderef{ExtensionBaseClaim45}; it has not been added as a hypothesis
of this lemma.

Suppose first that \(p>0\).  All other hypotheses of
\noderef{ExtensionBaseClaim45} are the facts just established: (5),
Case 2 and (19)--(23), (b1)--(b4), the exact support enumeration and
injective chain, the terminal and preterminal PE fields, and the
contrary component path.  Apply it to obtain its exact \(\varrho_3\),
including properness, unique uncoloredness, (d1), the complete prefix
agreements (d2) and (d2z), and (d3).

For completeness, if \(p=0\), the restoration required after
\(\mu_0\) has no preterminal slot and hence no Algorithm C step.  Set
\(\varrho_3=\mu_0\).  We check directly the three claimed clauses,
rather than pretending that a preterminal index exists.  Equation (19)
and the chosen membership \(\beta\in\overline\varrho_1(u)\) say that
\(\delta_n\) and \(\beta\), respectively, are missing on \(T'_j\), and
they are distinct because \(\delta_n\in S(L_1)\) whereas
\(\beta\notin S(L_1)\); here the passage from the two vertex witnesses to
tree-missing colors is exactly \noderef{TreeMissingColors}.  Thus
\noderef{OutsideTreeColorSwapStable} makes \(\varrho_0\) proper,
one-edge-uncolored, stable on \(T'_j\), and pointwise equal to
\(\varrho_1\) on every exact edge identity incident with \(T'_j\).

We identify the component used in the second swap.  In the rooted old
path \(P\), the internal vertex \(u\) has one path-edge identity toward
\(x\), namely the \(\gamma_n\)-colored exit identity supplied by
\noderef{ExtensionBaseCaseTwoObservation21}, and a distinct path-edge
identity toward \(v_n\).  Both are incident with \(u\), hence adjacent
by \noderef{AdjacentEdges}.  The second identity has one of the two old
path colors by \noderef{TwoColorIncidenceGraph}; since the first already
has color \(\gamma_n\), properness in
\noderef{IsProperEdgeColoring} forces the second identity to have color
\(\delta_n\).  The outside swap fixes both identities because \(u\in
V(T'_j)\).  Every other vertex-node of the exit suffix \(Q-u\) is outside
\(V(T_n)\cup V(T'_j)\).  Hence
\noderef{OutsideTreeColorSwap} fixes every \(\gamma_n\)-identity on that
suffix and changes precisely its outside \(\delta_n\)-identities to
\(\beta\).  Thus \(Q\) is a \((\gamma_n,\beta)\)-path under
\(\varrho_0\).  It cannot extend from \(u\) toward \(v_n\): the back
identity remains \(\delta_n\), while stability of the outside swap and
\noderef{MissingColors} keep \(\beta\) missing at \(u\); moreover
properness allows no second incident \(\gamma_n\)-identity.

There is no extension at the other endpoint \(x\) either.  Since
\(x\notin V(T'_j)\), every new \(\gamma_n\)- or \(\beta\)-identity
incident there corresponds under the outside transposition to an old
\(\gamma_n\)- or \(\delta_n\)-identity.  Its edge-node therefore belonged
to the full old component, whose support is exactly \(P\); the degree-one
endpoint calculation at \(x\) says that it is the terminal identity already
used by \(Q\).  If the new
\((\gamma_n,\beta)\)-component through \(u\) had a node outside \(Q\), a
shortest route from such a node to \(Q\) would first meet an internal node
of \(Q\), giving that node a third neighbor, or one of the endpoints,
extending it.  The internal case contradicts
\noderef{TwoColorIncidenceDegreeAtMostTwo}; the endpoint cases were just
excluded at \(u\) and \(x\).  By the component definition in
\noderef{TwoColorComponent}, the switched component \(K\) is therefore
exactly the support of \(Q\), and observation (21) gives
\[
 K=\{r:r\text{ is a node of }Q\},\qquad
 K\cap\{\mathsf{inl}(y):y\in V(T_n)\cup V(T'_j)\}
     =\{\mathsf{inl}(u)\}.                                    \tag{9}
\]

Since \(\gamma_n\ne\beta\), the properness conclusion of
\noderef{TwoColorComponentSwapProper} makes \(\varrho_3=\mu_0\) proper.
The pointwise definition \noderef{TwoColorComponentSwap} exchanges two
colored values and fixes the uncolored value, so it leaves exactly the
same single edge uncolored as \(\varrho_0\), in the sense of
\noderef{IsOneEdgeUncolored}.

We can now prove (d1) from (9).  Every vertex of
\(T_j(v_n)-v_n\) lies in \(V(T_n)\) and differs from \(u\), by (8).
The outside swap fixes all exact identities incident with that prefix.
If an edge-node of \(K\) were incident with one of its vertices, the
adjacency rule of \noderef{TwoColorIncidenceGraph} would put that
vertex-node in \(K\), contradicting (9).  The component swap therefore
fixes every incident identity as well.  The witness definition
\noderef{PresentColors} and complement definition
\noderef{MissingColors} turn these incident-edge equalities into equality
of the complete missing sets.  Both clauses of (b1) consequently transfer
from \(\varrho_1\) to \(\varrho_3\), proving (d1).

For (d2), first consider a vertex
\(y\in V(T_j(z))\setminus V(T'_j)\) with \(y\ne z\).  Observation (23)
makes \(\beta\) present at \(y\) under \(\varrho_1\).  Equation (2) says
that \(\delta_n\) is present throughout \(T_j\) under
\(\varphi_{j-1}\); since \(y\ne v_n\) (the seed containing \(v_n\)
lies in \(T'_j\)), (b2) transfers this presence to \(\varrho_1\).
Both \(\beta\) and \(\delta_n\) are missing on the closed tree \(T'_j\),
witnessed by \(u\) and \(v_n\).  Thus the closedness clause of
\noderef{TashkinovClosure}, interpreted by
\noderef{ColoredVertexSetClosed}, forbids either color on a boundary
identity of \(V(T'_j)\).  Every incident identity of either color at the
outside vertex \(y\) which entered \(V(T'_j)\) would be such a boundary
identity by \noderef{BoundaryEdges}; consequently every one is wholly
outside \(T'_j\).  The outside swap
transposes the two complete nonempty incidence sets at \(y\); it therefore
preserves the full present-color set, and hence the full missing-color set
by \noderef{PresentColors} and \noderef{MissingColors}.  If instead
\(y\in V(T'_j)\setminus\{u\}\), the first swap fixes every incident edge
and gives the same missing-set equality directly.  Finally (9) says that
the second swap fixes every identity incident with either kind of vertex.
This proves the d2 missing-set equality away from \(u,z\).

The exact displayed edges are fixed as well.  Every seed identity lies in
\(T'_j\), so the outside swap fixes it.  At an appended position of
\(T_j(z)\), observation (23) excludes \(\beta\) from the exact edge.
The original augmentation certificate and equation (2) exclude
\(\delta_n\) from its \(\varphi_{j-1}\)-color, and (b2) transfers that
exact value to \(\varrho_1\); hence the outside swap fixes this identity.
It also fixes every displayed identity of \(T'_j(u)\).  No displayed edge
of \(T_j(z)\) lies in \(K\), because both endpoints lie in \(V(T_n)\)
and adjacency would contradict (9).  No displayed edge of \(T'_j(u)\)
lies in \(K\) either: both endpoints lie in \(V(T'_j)\), so (9) would
force both to equal \(u\), contrary to looplessness.  Thus the second
swap also fixes every exact edge of both prefixes, completing the ordinary
d2 agreements.

It remains to check the two endpoint refinements.  Suppose first that
\(\beta\in\overline\varrho_1(z)\).  If \(z\in V(T'_j)\), then the
distinct vertices \(z\) and \(u\) of the elementary set \(V(T'_j)\)
would both miss \(\beta\): they are distinct because
\(z\in V(T_j)\subseteq V(T_n)\) while \(u\notin V(T_n)\).  By
\noderef{ColoredVertexSetElementary}, this membership case makes the
antecedent impossible.  If
\(z\notin V(T'_j)\), then \(z\ne v_n\), and equation (2) plus (b2) make
\(\delta_n\) present at \(z\) under \(\varrho_1\).  Closedness of
\(T'_j\), in the sense of \noderef{ColoredVertexSetClosed}, applied to
the \(\beta\)- and \(\delta_n\)-missing witnesses
\(u,v_n\), excludes both colors from identities between \(z\) and
\(T'_j\), which are boundary identities by \noderef{BoundaryEdges}.  The
outside swap therefore changes every incident
\(\delta_n\)-identity at \(z\) to \(\beta\), while there is no incident
\(\beta\)-identity to change back; by \noderef{PresentColors} and
\noderef{MissingColors}, \(\delta_n\) becomes missing at \(z\) under
\(\varrho_0\).  Since \(z\in V(T_n)\) and (9) contains no vertex-node of
\(T_n\), the second swap fixes every identity incident with \(z\).
Therefore \(\delta_n\in\overline\varrho_3(z)\), proving the conditional
part of (d2).

If \(z\) is displayed strictly before \(u\) in \(T'_j\), then
\(z\in V(T'_j)\setminus\{u\}\).  The first swap fixes every incident
identity there, and (9) shows that the second swap does too.  The same
\noderef{PresentColors}--\noderef{MissingColors} calculation gives
\(\overline\varrho_3(z)=\overline\varrho_1(z)\), which is (d2z).
Finally \(p=0\) makes \(\gamma_{1_1}=\gamma_n\).  At \(u\), the first
swap preserves the \(\beta\)-missing condition and the exact present
\(\gamma_n\)-exit identity.  Properness makes that identity the unique
incident \(\gamma_n\)-identity, and (9) puts it in the switched component.
The component swap changes it to \(\beta\).  The definitions
\noderef{PresentColors} and \noderef{MissingColors} therefore make
\(\gamma_{1_1}=\gamma_n\) missing at \(u\), which is (d3).  Thus in both
cases we
have a coloring \(\varrho_3\) with the exact output clauses used below.

Observation (23) places the displayed position \(q_z\) at or after the
terminal position of the seed \(T_j(v_n)\).  We split now according as
\(z\ne v_n\) or \(z=v_n\), before using any result that assumes strict
order.

Suppose first that \(z\ne v_n\).  The displayed vertices of
\noderef{TreeSequence} are injective, so the weak position inequality is
strict: \(v_n\) occurs strictly before \(z\).  If \(j>1\), use the
certified lower ETT \(B\) already returned by
\noderef{CertifiedETTRestriction} at target rung \(j-1\).  It has the
original input \(\varphi'_0=\varphi_0\), so its instance of assertion
(vi) is the one supplied to
\noderef{ExtensionBaseCaseTwoContradiction}.  Its assertion-(i) premise is
supplied differently, because that premise ranges over corresponding
certificates: for each such certificate \(C_3\), instantiate the input
argument with \(C_3.\mathrm{replacementInput}\) and the lower-ETT argument
with \(C_3.\mathrm{replacementETT}\).  The equality
\(C_3.\mathrm{rungAgreement}\) and the original strict stage bound for
\(B\) prove that this replacement ETT has fewer than \(n\) rungs; its ETT
certificate gives positivity, and the maximum-property argument passed by
the contradiction node is the one obtained from
\noderef{MaximumPropertyTransport}.  The broadened package therefore
extracts assertion (i) for exactly \(C_3.\mathrm{replacementETT}\), not for
an ETT tacitly assumed to retain the original input.  The PE versus SE/RE
constructor supplies the exact seed and connecting-edge data, while (b2),
(19), and the supporting-color transport supply the remaining prefix and
missing-set data.  Together with the strict order just proved, every
hypothesis of \noderef{ExtensionBaseCaseTwoContradiction} is now satisfied
by the actual \(\varrho_3\).  That node rebuilds the copied augmentations,
transports maximum property, appends the remaining exact \(T_j(z)\)-edges,
and gives a contradiction.

If \(z\ne v_n\) and \(j=1\), perform the same finite reconstruction from
the initial edge.  Clause (d1) copies each initial augmentation to
\(T_1(v_n)\); (d2) and (d2z) copy every retained closure augmentation
through \(u\).  The supporting-color inclusion and (d3) give
\[
 \overline\varphi_0(v_n)\subseteq
   \overline\varrho_3(v_n)\cup\overline\varrho_3(u).
\]
Scanning the remaining exact edges through \(T_1(z)\), exactly as in the
Case 1 scan, supplies a \(\varrho_3\)-augmentation certificate at each new
vertex.  Properness and unique uncoloredness of \(\varrho_3\), together
with the unchanged graph, edge, and numerical clauses of the original
\noderef{KTriple}, make \((G,e_0,\varrho_3)\) a \(k\)-triple.  The resulting
tree is a Tashkinov tree by
\noderef{TashkinovAugmentationSequencePreservesTashkinovTree}.  If \(z\)
is terminal, this tree contains \(T_1\) and \(u\notin V(T_1)\), contrary
to \(|V(T_1)|=\noderef{TashkinovOrder}(G)\).  Otherwise observation (23)
says that \(\beta\) is missing at \(z\); the conditional part of (d2)
makes \(\delta_n\) missing there, while the ordinary d2 equality at the
distinct earlier vertex \(v_n\), together with (19), makes
\(\delta_n\) missing at \(v_n\).  The two distinct vertices lie in the
constructed tree, contradicting its elementariness from
\noderef{TashTree}.

It remains to handle \(z=v_n\).  Observation (23) searches only positions
strictly after the seed for a first \(\beta\)-missing vertex; therefore
\(v_n\) can be selected only in its terminal alternative.  Display
injectivity from \noderef{TreeSequence} and the ordered equality in
\noderef{TreeSequenceInitialSegment} then force
\[
                 T_j(z)=T_j(v_n)=T_j.                         \tag{10}
\]
The vertex \(u\) lies in \(V(T'_j)\setminus V(T_n)\), whereas
\noderef{TreeSequenceInitialSegment} and
\noderef{TreeSequenceVertexSet} put the whole seed in \(V(T_n)\).  Since
the seed is an ordered initial segment of \(T'_j\), \(v_n=z\) is displayed
strictly before \(u\).  Clause (d2z)
therefore gives
\[
              \overline\varrho_3(v_n)=\overline\varrho_1(v_n), \tag{11}
\]
in addition to the d2 equality on every exact displayed edge through
\(T'_j(u)\).

We reconstruct that prefix under \(\varrho_3\) one augmentation at a
time.  Stop the \(\varrho_1\)-augmentation sequence defining the closure
\(T'_j\) at the step which creates \(u\).  At a retained step, its exact
new edge has the same value under \(\varrho_3\) by (d2).  Its old
missing-color witness occurs strictly before \(u\).  If that witness is
\(v_n=z\), equation (11) preserves its full missing set; otherwise it is
neither \(u\) nor \(z\), and the ordinary d2 equality preserves its full
missing set.  Each exact step is consequently a
\(\varrho_3\)-\noderef{TashkinovAugmentation}, and finite induction gives
a \(\varrho_3\)-augmentation sequence from \(T_j=T_j(v_n)\) through the
prefix \(T'_j(u)\).

When \(j>1\), d1 restricts to
\((T_{j-1},D_{j-1},\varphi_{j-1})\)-stability.  At every original
augmentation from the prescribed PE seed, or from the prescribed SE/RE
one-edge seed, to \(T_j(v_n)\), its old missing-color witness lies before
\(v_n\); d1 preserves that witness and the exact new edge value.  Thus the
same stage route is valid under \(\varrho_3\).  Extend \(T'_j(u)\) to a
\(\varrho_3\)-closure using \noderef{TashkinovClosureExists}.  The exact
augmentation sequences just obtained concatenate to make this a closure
of the exact PE seed, or of the exact SE/RE one-edge seed, prescribed at
stage \(j-1\).  The exact stage clause of the original
\noderef{MaximumProperty}, applied to the d1 restriction and this
reconstructed route, bounds the closure by
\(|V(T_j)|\).  But (10) puts all of \(V(T_j)\) in it, and it also contains
\(u\notin V(T_j)\).  Its order is therefore strictly larger than
\(|V(T_j)|\), a contradiction.

When \(j=1\), d1 similarly copies the initial augmentations through
\(T_1(v_n)=T_1\), and the reconstruction above appends the exact closure
prefix through \(u\).  Properness and unique uncoloredness of
\(\varrho_3\), together with the graph and numerical clauses of the
original \noderef{KTriple}, make \((G,e_0,\varrho_3)\) another
\(k\)-triple.  By
\noderef{TashkinovAugmentationSequencePreservesTashkinovTree} the result is
a Tashkinov tree under the proper one-edge-uncolored coloring
\(\varrho_3\).  It contains \(T_1\) and the additional vertex
\(u\notin V(T_1)\), contradicting the initial maximum equality
\(|V(T_1)|=\noderef{TashkinovOrder}(G)\).  Thus the equality case is also
impossible, and Case 2 is excluded for both \(j>1\) and \(j=1\).

Both alternatives in (6) are impossible.  Hence no
\(w\ne v_n\) in \(T_n\) belongs to the terminal
\((\gamma_n,\delta_n)\)-component.  Together with the root observation
at the start, this proves the required equivalence for every
\(w\in V(T_n)\).
\end{proof}
